\documentclass[11pt,a4paper]{article}
\usepackage[utf8]{inputenc}
\usepackage[T1]{fontenc}
\usepackage{amsmath,amssymb}
\usepackage{graphicx}
\usepackage{xcolor}
\usepackage{listings}
\usepackage{hyperref}
\usepackage{geometry}
\usepackage{booktabs}
\usepackage{fancyhdr}
\usepackage{tcolorbox}
\usepackage{algorithm}
\usepackage{algpseudocode}

\geometry{margin=1in}
\pagestyle{fancy}
\fancyhf{}
\rhead{\thepage}
\lhead{GraSMoS User Manual}

\hypersetup{
    colorlinks=true,
    linkcolor=blue,
    filecolor=magenta,      
    urlcolor=cyan,
    citecolor=green,
    pdftitle={GraSMoS User Manual},
    pdfauthor={GraSMoS Development Team},
}

\lstdefinestyle{jsonStyle}{
    basicstyle=\ttfamily\small,
    breaklines=true,
    backgroundcolor=\color{gray!10},
    frame=single,
    numbers=left,
    numberstyle=\tiny\color{gray},
    keywordstyle=\color{blue},
    stringstyle=\color{red},
    commentstyle=\color{green!50!black},
    morestring=[b]",
}

\lstdefinestyle{bashStyle}{
    basicstyle=\ttfamily\small,
    breaklines=true,
    backgroundcolor=\color{blue!5},
    frame=single,
    language=bash,
}

\title{\textbf{GraSMoS User Manual}\\
\large Graph-Aware Structure Search with Monte Carlo Simulation\\
Version 2.0}
\author{lipai@mail.sim.ac.cn}
\date{\today}

\begin{document}

\maketitle
\tableofcontents
\newpage

%==============================================================================
\section{Introduction}
%==============================================================================

\subsection{Overview}
GraSMoS (Graph-Aware Structure Search with Monte Carlo Simulation) is a sophisticated computational tool designed to find stable atomic structures through advanced global optimization algorithms. It combines Monte Carlo sampling with a climbing algorithm and graph-theoretic collective moves to efficiently explore the potential energy surface and identify low-energy structural configurations. Key innovations include bond-switching for topological rearrangement, community detection on bond graphs for coordinated multi-atom moves, and Laplacian soft-mode analysis for identifying low-energy deformation pathways.

\subsection{Key Features}
\begin{itemize}
    \item \textbf{Multi-calculator support}: Compatible with various energy calculators including EAM, CHGNet, DeepMD, FAIRChem, VASP, LAMMPS, and custom Python calculators
    \item \textbf{Graph-theoretic collective moves}: Bond-switching (generalized Stone-Wales), community detection (Louvain), Laplacian soft modes for efficient PES exploration
    \item \textbf{Adaptive mode weighting}: Multi-armed bandit optimization of move-type probabilities based on historical performance
    \item \textbf{Flexible mobility control}: Supports index-based and region-based constraints on atomic movement
    \item \textbf{Permutation-invariant structure comparison}: Uses sorted SOAP descriptors combined with energy gating for robust duplicate detection
    \item \textbf{Automatic geometry classification}: Detects cluster, wire, slab, or bulk structures with appropriate handling of vacuum axes
    \item \textbf{Energy-weighted sampling}: Prioritizes movement of unstable atoms through per-atom energy analysis
\end{itemize}

%==============================================================================
\section{Theoretical Background}
%==============================================================================

\subsection{Stochastic Surface Walking (SSW) Algorithm}
GraSMoS is based on the Stochastic Surface Walking method \cite{shang2013ssw}, which combines:

\begin{enumerate}
    \item \textbf{Random displacement}: Structure perturbation along a weighted random direction
    \item \textbf{Climbing phase}: Addition of Gaussian bias potentials to escape local minima
    \item \textbf{Local optimization}: Minimization to nearest stationary point
    \item \textbf{Metropolis acceptance}: Temperature-dependent Monte Carlo acceptance
\end{enumerate}

\subsection{Algorithm Workflow}

\begin{algorithm}
\caption{GraSMoS Search Algorithm}
\begin{algorithmic}[1]
\State Initialize structure $\mathbf{R}_0$ and energy pool $E = \emptyset$
\State Optimize $\mathbf{R}_0 \rightarrow \mathbf{R}_{\text{min}}$, add to pool
\For{$\text{step} = 1$ to $N_{\text{MC}}$}
    \State Generate random direction $\mathbf{N}$ (energy-weighted)
    \State Displace: $\mathbf{R}' = \mathbf{R}_{\text{current}} + d_s \mathbf{N}$
    \For{$h = 1$ to $H$ (Gaussian potentials)}
        \State Add Gaussian bias: $V_{\text{bias}} = \sum_{i=1}^{h} w e^{-|\mathbf{R} - \mathbf{R}_i|^2 / (2\sigma^2)}$
        \State Optimize on biased PES: $\mathbf{R}_h \rightarrow \mathbf{R}_{\text{climb}}$
        \If{converged or max steps reached}
            \State \textbf{break}
        \EndIf
    \EndFor
    \State Optimize on real PES: $\mathbf{R}_{\text{climb}} \rightarrow \mathbf{R}_{\text{new}}$
    \State Compute $E_{\text{new}} = E(\mathbf{R}_{\text{new}})$
    \If{$\mathbf{R}_{\text{new}}$ is not duplicate}
        \State Add $(\mathbf{R}_{\text{new}}, E_{\text{new}})$ to pool
    \EndIf
    \State Accept/reject via Metropolis: $P = \min(1, e^{-\Delta E / k_B T})$
    \State Update $\mathbf{R}_{\text{current}}$
\EndFor
\State \Return lowest energy structure from pool
\end{algorithmic}
\end{algorithm}

\subsection{Energy-Weighted Direction Sampling}
Unlike traditional Maxwell-Boltzmann sampling, GraSMoS uses per-atom potential energies to guide sampling:

\begin{center}
\textbf{N}_i \ \sim\  N(0,\ \sigma_i^2),\qquad \sigma_i = \sqrt{k_B T} \cdot \exp\!\bigl(E_i^{\,\text{norm}}\bigr)
\end{center}

where $E_i^{\text{norm}}$ is the normalized energy of atom $i$ relative to the lowest-energy atom of the same element. This prioritizes movement of high-energy, unstable atoms.

\subsection{Structure Comparison via SOAP Descriptors}
To handle atomic permutations, GraSMoS uses permutation-invariant descriptors:

\begin{enumerate}
    \item Compute per-atom SOAP vectors $\{\mathbf{s}_i\}_{i=1}^N$
    \item Sort by L2 norm: $\mathbf{s}_{\sigma(1)}, \ldots, \mathbf{s}_{\sigma(N)}$ where $\|\mathbf{s}_{\sigma(i)}\| \leq \|\mathbf{s}_{\sigma(i+1)}\|$
    \item Flatten to structure descriptor: $\mathbf{D} = [\mathbf{s}_{\sigma(1)}; \ldots; \mathbf{s}_{\sigma(N)}]$
    \item Compare: structures are duplicates if $\|\mathbf{D}_1 - \mathbf{D}_2\| < \tau_{\text{desc}}$ and $|E_1 - E_2| < \tau_E$
\end{enumerate}

%==============================================================================
\section{Installation}
%==============================================================================

\subsection{System Requirements}
\begin{itemize}
    \item Python 3.8 or higher
    \item pip (Python package manager)
    \item ASE $\geq$ 3.26.0 (Atomic Simulation Environment)
    \item dscribe $\geq$ 2.1.0 (for SOAP descriptors)
\end{itemize}

\subsection{Installation Steps}

\begin{lstlisting}[style=bashStyle, caption={Installation from source}]
# Clone the repository
git clone https://github.com/lipai-ustc/GraSMoS.git
cd grasmos

# (Optional) Create Conda environment
conda create -n grasmos python=3.10 -y
conda activate grasmos

# Install GraSMoS
pip install .

# (Optional) Install additional calculators
pip install deepmd-kit fairchem-core
\end{lstlisting}

\begin{tcolorbox}[colback=yellow!10,colframe=orange!50,title=Development Mode]
For active development, use editable installation:
\begin{lstlisting}[style=bashStyle]
pip install -e .
\end{lstlisting}
This allows code modifications without reinstallation.
\end{tcolorbox}

%==============================================================================
\section{Basic Usage}
%==============================================================================

\subsection{Input Files}
GraSMoS requires two input files in the working directory:

\begin{enumerate}
    \item \texttt{input.json}: Calculation parameters configuration
    \item \texttt{init.xyz}: Initial atomic structure in XYZ format (default path, can be changed in configuration)
\end{enumerate}

\subsection{Running GraSMoS}
\begin{lstlisting}[style=bashStyle, caption={Basic execution}]
# Navigate to directory containing input.json and init.xyz
cd /path/to/working/directory

# Run search
grasmos
\end{lstlisting}

\subsection{Output Configuration}

The \texttt{output} section controls output settings:

\begin{table}[h]
\centering
\begin{tabular}{lll}
\toprule
\textbf{Parameter} & \textbf{Description} & \textbf{Default} \\
\midrule
\texttt{directory} & Output directory name & "grasmos\_output" \\
\texttt{rd\_xyz} & Enable output of random direction information & false \\
\texttt{debug} & Enable debug mode for detailed logging & false \\
\bottomrule
\end{tabular}
\caption{Output configuration parameters}
\end{table}

\subsection{Output Files}
Results are saved in the specified output directory:

\begin{itemize}
    \item \texttt{all\_minima.xyz}: All discovered minimum energy structures (trajectory format)
    \item \texttt{best\_str.xyz}: Lowest energy structure
    \item \texttt{grasmos\_log.txt}: Detailed search log with energies and algorithm status
    \item \texttt{rd\_info.xyz}: Random direction information (only if \texttt{rd\_xyz} is true)
\end{itemize}

\subsection{Debug Mode Output}

In debug mode, the log file contains detailed information for each optimization step:

\begin{lstlisting}
Step 1: E_total = -125.234567 eV, E_base = -125.500000 eV, E_bias = 0.250000 eV, E_wall = 0.015433 eV, fmax = 0.850000 eV/\AA
Step 2: E_total = -125.456789 eV, E_base = -125.650000 eV, E_bias = 0.180000 eV, E_wall = 0.013211 eV, fmax = 0.420000 eV/\AA
...
\end{lstlisting}

%==============================================================================
\section{Configuration Reference}
%==============================================================================

\subsection{input.json Structure}
The configuration file is divided into several logical sections:

\begin{lstlisting}[style=jsonStyle, caption={Minimal input.json example}]
{
  "system": {
    "name": "MySystem",
    "task": "global_search",
    "structure": "init.xyz"
  },
  "potential": {
    "type": "eam",
    "model": "potential.eam.alloy"
  },
  "monte_carlo": {
    "steps": 100,
    "temperature": 300
  },
  "climbing": {
    "gaussian": {
      "height": 0.2,
      "width": 0.2,
      "Nmax": 20
    },
    "random_direction": {
      "mode": ["thermo", "atomic"],
      "ratio": [[[0.5, 0.5], 1]]
    }
  },
  "optimizer": {
    "max_steps": 500,
    "fmax": 0.05
  },
  "output": {
    "directory": "grasmos_output",
    "debug": false
  }
}
\end{lstlisting}

\subsection{Potential Configuration}

\subsubsection{EAM Potential}
\begin{lstlisting}[style=jsonStyle]
"potential": {
  "type": "eam",
  "model": "AlCu.eam.alloy"
}
\end{lstlisting}

\subsubsection{CHGNet}
\begin{lstlisting}[style=jsonStyle]
"potential": {
  "type": "chgnet",
  "model": "pretrained"
}
\end{lstlisting}

\subsubsection{DeepMD}
\begin{lstlisting}[style=jsonStyle]
"potential": {
  "type": "deepmd",
  "model": "dp_model.pb"
}
\end{lstlisting}

\subsubsection{FAIRChem (Open Catalyst Project)}
\begin{lstlisting}[style=jsonStyle]
"potential": {
  "type": "fairchem",
  "model": "EquiformerV2-31M-S2EF-OC20-All+MD",
  "device": "cuda",
  "task_name": "oc20"
}
\end{lstlisting}

\subsubsection{VASP}
\begin{lstlisting}[style=jsonStyle]
"potential": {
  "type": "vasp",
  "model": "INCAR"
}
\end{lstlisting}
Reads VASP parameters from the specified INCAR file.

\subsubsection{Custom Python Calculator}
\begin{lstlisting}[style=jsonStyle]
"potential": {
  "type": "python"
}
\end{lstlisting}
Requires a \texttt{calculator.py} file in the working directory defining a \texttt{calculator} variable with an ASE Calculator object.

\subsection{Monte Carlo Parameters}
\begin{table}[h]
\centering
\begin{tabular}{lll}
\toprule
\textbf{Parameter} & \textbf{Description} & \textbf{Default} \\
\midrule
\texttt{steps} & Total MC steps & Required \\
\texttt{temperature} & Temperature (K) for Metropolis & Required \\
\bottomrule
\end{tabular}
\caption{Monte Carlo configuration parameters}
\end{table}

\subsection{Climbing Phase Parameters}

\subsubsection{Gaussian Bias Potentials}
\begin{table}[h]
\centering
\begin{tabular}{lll}
\toprule
\textbf{Parameter} & \textbf{Description} & \textbf{Default} \\
\midrule
\texttt{gaussian.height} & Height of Gaussian bias potentials (w parameter) in eV & 0.2 \\
\texttt{gaussian.width} & Width of Gaussian potentials (ds parameter) in \AA & 0.2 \\
\texttt{gaussian.Nmax} & Maximum number of Gaussians per climbing (H parameter) & 20 \\
\bottomrule
\end{tabular}
\caption{Gaussian bias potential parameters}
\end{table}

\subsubsection{Random Direction Generation}
\begin{table}[h]
\centering
\begin{tabular}{lll}
\toprule
\textbf{Parameter} & \textbf{Description} & \textbf{Default} \\
\midrule
\texttt{mode} & List of methods for generating random search directions & ["thermo", "atomic"] \\
\texttt{ratio} & Weight configurations for combining multiple modes & [[[0.5, 0.5], 1]] \\
\texttt{rotation\_param} & Dimer rotation parameter & 10 \\
\texttt{element\_weights} & Element-specific weight factors & {} \\
\texttt{adaptive} & Enable adaptive mode weighting (v2.0+) & false \\
\texttt{adaptive\_interval} & Steps between weight recomputation & 10 \\
\texttt{adaptive\_alpha} & Acceptance-rate exponent for scoring & 2.0 \\
\texttt{adaptive\_floor} & Minimum exploration weight fraction & 0.05 \\
\texttt{adaptive\_smoothing} & EMA smoothing factor & 0.7 \\
\bottomrule
\end{tabular}
\caption{Random direction generation parameters}
\end{table}

\subsection{Graph-Theoretic Collective Moves (v2.0)}
GraSMoS introduces graph-based methods for identifying and executing effective multi-atom cooperative motions:

\begin{itemize}
    \item \textbf{Bond Graph Construction}: A weighted adjacency matrix is built from covalent radii, with edge weight $w = e^{-d / 0.5 r_{\text{cov}}}$. Edges with $w < 0.01$ are dropped. Per-atom energies (when available) modulate edge weights, biasing analysis toward high-energy regions.

    \item \textbf{Community Detection (\texttt{community} mode)}: Louvain community detection on the bond graph identifies natural clusters (functional groups, coordination polyhedra, ring systems). A randomly selected community undergoes rigid translation, rotation, or radial breathing.

    \item \textbf{Laplacian Soft Modes (\texttt{laplacian} mode)}: Low-frequency eigenvectors of the normalized graph Laplacian $L = I - D^{-1/2}AD^{-1/2}$ identify collective deformation patterns that cost minimal energy in a harmonic approximation.

    \item \textbf{Bond Switching (\texttt{bond\_switch} mode)}: Large-angle (30°--90°) relative rotation of two local fragments about a bond, with graph-based candidate screening (coordination $\ge 2$, each endpoint has at least one other neighbour). This is the material-agnostic generalization of the Stone-Wales transformation.

    \item \textbf{Adaptive Mode Weighting}: A multi-armed bandit approach dynamically adjusts per-mode sampling probabilities based on acceptance rate, energy improvement, and new-minima discovery. Modes are never eliminated (exploration floor prevents this), and EMA smoothing prevents abrupt drift.
\end{itemize}
\caption{Random direction generation parameters}
\end{table}

\subsubsection{Random Direction Modes}
\begin{table}[h]
\centering
\begin{tabular}{ll}
\toprule
\textbf{Mode} & \textbf{Description} \\
\midrule
\texttt{thermo} & Temperature-based random vectors (Boltzmann distribution) \\
\texttt{atomic} & Energy-weighted thermo; high-energy atoms get larger displacements \\
\texttt{nl} & Non-local pair attraction for bond formation \\
\texttt{atnl} & Atomic-energy guided non-local pair \\
\texttt{bond\_rotation} & Coordinated bond rotation (generalized GSW-like) \\
\texttt{bond\_switch} & Generalized bond-switching (Stone-Wales-like), large-angle topology change \\
\texttt{shell} & Coordination shell collective motion (breathing/shear) \\
\texttt{community} & Louvain community detection on bond graph; rigid-body cluster moves \\
\texttt{laplacian} & Graph Laplacian soft-mode deformation \\
\texttt{python} & Custom user-defined function \\
\bottomrule
\end{tabular}
\caption{Random direction generation modes}
\end{table}

\subsection{Mobile Control}

Mobile control allows constraining which atoms can move during optimization. Four modes are supported:

\subsubsection{Mode 1: All Atoms Mobile (Default)}
\begin{lstlisting}[style=jsonStyle]
"mobile_control": {
  "mode": "all"
}
\end{lstlisting}
or omit the \texttt{mobile\_control} section entirely.

\subsubsection{Mode 2a: Specify Mobile Atoms}
\begin{lstlisting}[style=jsonStyle]
"mobile_control": {
  "mode": "indices_free",
  "indices_free": [10, 20, 25],
  "wall_strength": 10.0,
  "wall_offset": 2.0
}
\end{lstlisting}
Only listed atoms (0-based indexing) are mobile; all others are fixed.

\subsubsection{Mode 2b: Specify Fixed Atoms}
\begin{lstlisting}[style=jsonStyle]
"mobile_control": {
  "mode": "indices_fix",
  "indices_fix": [0, 1, 2],
  "wall_strength": 10.0,
  "wall_offset": 2.0
}
\end{lstlisting}
Listed atoms are fixed; all others are mobile.

\subsubsection{Mode 3: Region-Based Control}

\paragraph{Spherical Region}
\begin{lstlisting}[style=jsonStyle]
"mobile_control": {
  "mode": "region",
  "region_type": "sphere",
  "center": [5.0, 5.0, 5.0],
  "radius": 10.0,
  "wall_strength": 10.0,
  "wall_offset": 2.0
}
\end{lstlisting}
Atoms within the sphere are mobile.

\paragraph{Slab Region}
\begin{lstlisting}[style=jsonStyle]
"mobile_control": {
  "mode": "region",
  "region_type": "slab",
  "origin": [0.0, 0.0, 0.0],
  "normal": [0, 0, 1],
  "min_dist": -5.0,
  "max_dist": 5.0,
  "wall_strength": 10.0,
  "wall_offset": 2.0
}
\end{lstlisting}
Atoms between two parallel planes are mobile.

\paragraph{Lower Region}
\begin{lstlisting}[style=jsonStyle]
"mobile_control": {
  "mode": "region",
  "region_type": "lower",
  "axis": "z",
  "threshold": 10.0,
  "wall_strength": 10.0,
  "wall_offset": 2.0
}
\end{lstlisting}
Atoms with coordinate $\le$ threshold along the specified axis are mobile.

\paragraph{Upper Region}
\begin{lstlisting}[style=jsonStyle]
"mobile_control": {
  "mode": "region",
  "region_type": "upper",
  "axis": "z",
  "threshold": 10.0,
  "wall_strength": 10.0,
  "wall_offset": 2.0
}
\end{lstlisting}
Atoms with coordinate $\ge$ threshold along the specified axis are mobile.

\subsubsection{Wall Potential}
When \texttt{wall\_strength} $> 0$, a quadratic repulsive potential prevents mobile atoms from penetrating into immobile regions:

\begin{center}
\textbf{Wall potential:}\quad
$V_{\text{wall}} = 0$ \ if\ $d \leq d_{\text{offset}}$;\qquad
$V_{\text{wall}} = \tfrac{1}{2}k_{\text{wall}}(d - d_{\text{offset}})^2$ \ if\ $d > d_{\text{offset}}$
\end{center}

where $d$ is the penetration distance, $d_{\text{offset}}$ is \texttt{wall\_offset}, and $k_{\text{wall}}$ is \texttt{wall\_strength}.

%==============================================================================
\section{Low-Dimensional Systems}
%==============================================================================

\subsection{Geometry Classification}
GraSMoS automatically classifies structures based on vacuum analysis:
\begin{itemize}
    \item \textbf{Cluster (0D)}: Vacuum along all three axes
    \item \textbf{Wire (1D)}: Vacuum along two axes
    \item \textbf{Slab (2D)}: Vacuum along one axis
    \item \textbf{Bulk (3D)}: No significant vacuum
\end{itemize}

\subsection{Structure Preparation}
\begin{tcolorbox}[colback=red!5,colframe=red!50,title=IMPORTANT: Structure Centering]
For low-dimensional systems, you \textbf{must manually center} the structure in the simulation box before running GraSMoS. The geometry detection algorithm measures distances from atoms to cell boundaries; misplaced structures will be misclassified.

\textbf{Recommended placement:}
\begin{itemize}
    \item \textbf{0D cluster}: Center at $(L_x/2, L_y/2, L_z/2)$
    \item \textbf{1D wire}: Wire extends along one axis, centered in the other two
    \item \textbf{2D slab}: Centered along the vacuum direction
\end{itemize}
\end{tcolorbox}

\subsection{Vacuum Axis Handling}
For structures with vacuum axes, GraSMoS automatically:
\begin{enumerate}
    \item Detects vacuum axes using absolute coordinate analysis
    \item Centers the structure along vacuum axes to the box center
    \item Applies recentering after each accepted Monte Carlo move
    \item Prevents translational drift during search
\end{enumerate}

%==============================================================================
\section{Examples}
%==============================================================================

\subsection{Example 1: AlCu Cluster with EAM}
\begin{lstlisting}[style=bashStyle]
cd examples/AlCu-EAM
python generate_structure.py  # Generate init.xyz
grasmos                         # Run search
\end{lstlisting}

Configuration highlights:
\begin{itemize}
    \item EAM potential for Al-Cu system
    \item 100 Monte Carlo steps
    \item Temperature: 300 K
\end{itemize}

\subsection{Example 2: C$_{60}$ with DeepMD}
\begin{lstlisting}[style=bashStyle]
cd examples/C60-deepmd
# Ensure dp_model.pb is present
grasmos
\end{lstlisting}

Configuration highlights:
\begin{itemize}
    \item DeepMD neural network potential
    \item Requires \texttt{dp\_model.pb} trained model
    \item Install: \texttt{pip install deepmd-kit}
\end{itemize}

\subsection{Example 3: Au Surface with FAIRChem (Custom Calculator)}
\begin{lstlisting}[style=bashStyle]
cd examples/Au100-python-fairchem
# Check calculator.py for calculator definition
grasmos
\end{lstlisting}

Configuration highlights:
\begin{itemize}
    \item Custom Python calculator using FAIRChem
    \item Slab region mobility control for surface atoms
    \item Requires: \texttt{pip install fairchem-core}
\end{itemize}

%==============================================================================
\section{Advanced Topics}
%==============================================================================

\subsection{Task Types}
GraSMoS supports two task types:

\begin{itemize}
    \item \textbf{Global Search} (\texttt{global\_search}): Finds stable structures across the potential energy surface
    \item \textbf{Structure Sampling} (\texttt{structure\_sampling}): Samples atomic structures around existing minima
\end{itemize}

\subsection{Custom Calculator Development}
To use a custom ASE-compatible calculator:

\begin{enumerate}
    \item Create \texttt{calculator.py} in your working directory
    \item Define a \texttt{calculator} variable
    \item Set \texttt{"type": "python"} in \texttt{input.json}
\end{enumerate}

Example \texttt{calculator.py}:
\begin{lstlisting}[language=Python, style=jsonStyle]
from ase.calculators.emt import EMT

# Define your calculator
calculator = EMT()
\end{lstlisting}

\subsection{Custom Random Direction Generation}

When \texttt{random\_direction.mode} is set to \texttt{"python"}, create a \texttt{generate\_random\_direction.py} file in your working directory:

\begin{lstlisting}[language=Python, style=jsonStyle]
import numpy as np
from ase import Atoms

def generate_random_direction(atoms: Atoms) -> np.ndarray:
    """
    Generate custom random direction vector for GraSMoS algorithm.
    
    Parameters:
        atoms: ASE Atoms object representing current structure
    
    Returns:
        N: 1D numpy array of size (3*n_atoms,) representing the direction vector
    """
    n_atoms = len(atoms)
    N = np.random.randn(3 * n_atoms)  # Your custom logic here
    return N
\end{lstlisting}

Then configure in \texttt{input.json}:

\begin{lstlisting}[style=jsonStyle]
"climbing": {
  "random_direction": {
    "mode": ["python"]
  }
}
\end{lstlisting}

\subsection{Parallel Execution}
GraSMoS is designed for serial execution per search. For parallel exploration:
\begin{itemize}
    \item Run multiple independent searches with different initial structures
    \item Use different random seeds (implicitly different due to system time)
    \item Combine results post-processing
\end{itemize}

\subsection{Troubleshooting}

\subsubsection{Structure Misclassification}
\textbf{Symptom:} Cluster treated as bulk or vice versa.

\textbf{Solution:} Ensure structure is centered in the simulation box. Check with:
\begin{lstlisting}[language=Python, style=jsonStyle]
from ase.io import read
atoms = read('init.xyz')
pos = atoms.get_positions()
center = pos.mean(axis=0)
cell = atoms.get_cell()
box_center = (cell[0] + cell[1] + cell[2]) / 2
print(f"Structure center: {center}")
print(f"Box center: {box_center}")
\end{lstlisting}

\subsubsection{Slow Convergence}
\textbf{Possible causes:}
\begin{itemize}
    \item \texttt{gaussian\_height} too small: increase to 0.2--0.5 eV
    \item \texttt{temperature} too low: try 500--1000 K
    \item \texttt{max\_gaussians} too small: increase to 20--30
\end{itemize}

\subsubsection{Too Many Duplicates}
\textbf{Possible causes:}
\begin{itemize}
    \item \texttt{duplicate\_tol} too loose: default is 0.01, try 0.005
    \item Energy tolerance too large: adjust in source (default 1e-3 eV)
\end{itemize}

%==============================================================================
\section{Parameter Tuning Guidelines}
%==============================================================================

\begin{table}[h]
\centering
\begin{tabular}{lll}
\toprule
\textbf{Parameter} & \textbf{Effect} & \textbf{Tuning Strategy} \\
\midrule
\texttt{gaussian\_height} & Barrier crossing ability & Increase for rugged PES \\
\texttt{gaussian\_width} & Step size & Smaller for precise, larger for exploration \\
\texttt{max\_gaussians} & Climbing persistence & Increase if stuck in shallow minima \\
\texttt{temperature} & Acceptance probability & Higher for aggressive search \\
\texttt{monte\_carlo.steps} & Total exploration & More steps = better coverage \\
\bottomrule
\end{tabular}
\caption{Parameter tuning guidelines}
\end{table}

%==============================================================================
\section{References}
%==============================================================================

\begin{thebibliography}{99}

\bibitem{shang2013ssw}
Shang, C., \& Liu, Z.-P. (2013).
\textit{Stochastic Surface Walking Method for Structure Prediction and Pathway Searching}.
Journal of Chemical Theory and Computation, 9(5), 1838--1845.

\bibitem{shang2013jcp}
Shang, C., \& Liu, Z.-P. (2013).
\textit{Stochastic surface walking method for global optimization of atomic clusters and biomolecules}.
The Journal of Chemical Physics, 139(24), 244104.

\bibitem{jctc2012}
Zhang, X.-J., \& Shang, C., \& Liu, Z.-P. (2012).
\textit{Double-Ended Surface Walking Method for Pathway Building and Transition State Location of Complex Reactions}.
Journal of Chemical Theory and Computation, 9(12), 5745--5753.

% ── Graph theory & structure search ──

\bibitem{blondel2008louvain}
Blondel, V.~D., Guillaume, J.-L., Lambiotte, R., \& Lefebvre, E. (2008).
\textit{Fast unfolding of communities in large networks}.
Journal of Statistical Mechanics: Theory and Experiment, 2008(10), P10008.

\bibitem{chung1997spectral}
Chung, F.~R.~K. (1997).
\textit{Spectral Graph Theory}.
CBMS Regional Conference Series in Mathematics, Vol.~92. American Mathematical Society.

\bibitem{jacobs1998pebble}
Jacobs, D.~J., \& Thorpe, M.~F. (1998).
\textit{Generic rigidity percolation: The pebble game}.
Physical Review Letters, 75(22), 4051--4054.

\bibitem{stone1986stonewales}
Stone, A.~J., \& Wales, D.~J. (1986).
\textit{Theoretical studies of icosahedral C$_{60}$ and some related species}.
Chemical Physics Letters, 128(5--6), 501--503.

\bibitem{newman2004networks}
Newman, M.~E.~J. (2004).
\textit{Detecting community structure in networks}.
The European Physical Journal B, 38(2), 321--330.

\bibitem{belkin2003laplacian}
Belkin, M., \& Niyogi, P. (2003).
\textit{Laplacian eigenmaps for dimensionality reduction and data representation}.
Neural Computation, 15(6), 1373--1396.

% ── Machine learning potentials ──

\bibitem{behler2007nnp}
Behler, J., \& Parrinello, M. (2007).
\textit{Generalized neural-network representation of high-dimensional potential-energy surfaces}.
Physical Review Letters, 98(14), 146401.

\bibitem{batzner2022e3nn}
Batzner, S., Musaelian, A., Sun, L., Geiger, M., Mailoa, J.~P., Kornbluth, M., Molinari, N., Smidt, T.~E., \& Kozinsky, B. (2022).
\textit{E(3)-equivariant graph neural networks for data-efficient and accurate interatomic potentials}.
Nature Communications, 13, 2453.

\bibitem{zhang2018deepmd}
Zhang, L., Han, J., Wang, H., Car, R., \& E, W. (2018).
\textit{Deep potential molecular dynamics: a scalable model with the accuracy of quantum mechanics}.
Physical Review Letters, 120(14), 143001.

\bibitem{schutt2018schnet}
Sch{\"u}tt, K.~T., Sauceda, H.~E., Kindermans, P.-J., Tkatchenko, A., \& M{\"u}ller, K.-R. (2018).
\textit{SchNet -- A deep learning architecture for molecules and materials}.
The Journal of Chemical Physics, 148(24), 241722.

\bibitem{gasteiger2021gemnet}
Gasteiger, J., Becker, F., \& G{\"u}nnemann, S. (2021).
\textit{GemNet: Universal directional graph neural networks for molecules}.
Advances in Neural Information Processing Systems, 34, 6790--6802.

\bibitem{merchant2023scaling}
Merchant, A., Batzner, S., Schoenholz, S.~S., Aykol, M., Cheon, G., \& Cubuk, E.~D. (2023).
\textit{Scaling deep learning for materials discovery}.
Nature, 624(7990), 80--85.

\end{thebibliography}

%==============================================================================
\section{License}
%==============================================================================

This project is licensed under the GNU General Public License v3.0 (GPL-3.0).

You may copy, distribute, and modify this software under the terms of the GPL-3.0.
For the full license text, see: \url{https://www.gnu.org/licenses/gpl-3.0.html}

%==============================================================================
\appendix
\section{Complete Configuration Example}
%==============================================================================

\begin{lstlisting}[style=jsonStyle, caption={Full input.json with all options}]
{
  "system": {
    "name": "CompleteExample",
    "task": "global_search",
    "structure": "init.xyz"
  },
  "potential": {
    "type": "deepmd",
    "model": "dp_model.pb",
    "device": "cpu"
  },
  "monte_carlo": {
    "steps": 200,
    "temperature": 500
  },
  "climbing": {
    "gaussian": {
      "height": 0.2,
      "width": 0.2,
      "Nmax": 20
    },
    "random_direction": {
      "mode": ["thermo", "atomic"],
      "ratio": [[[0.5, 0.5], 1]],
      "element_weights": {"Cu": 1.5, "Al": 1.0},
      "rotation_param": 10
    }
  },
  "optimizer": {
    "max_steps": 500,
    "fmax": 0.03
  },
  "mobile_control": {
    "mode": "region",
    "region_type": "sphere",
    "center": "center",
    "radius": 8.0,
    "wall_strength": 15.0,
    "wall_offset": 1.5
  },
  "output": {
    "directory": "my_output",
    "rd_xyz": false,
    "debug": true
  }
}
\end{lstlisting}

\end{document}
